\setcounter{page}{68}
\setcounter{chapter}{7}

\chapter{Movimiento circular uniforme / Velocidad relativa / Gravitaci\'on}
\label{cap:cap7}
\vspace{-14ex}

\raggedright
\section{Movimiento circular uniforme (MCU)}

El Movimiento Circular Uniforme (MCU) es un tipo de movimiento en el que un objeto describe una trayectoria circular con rapidez constante. Aunque la magnitud de la velocidad no cambia, la dirección del movimiento sí lo hace constantemente, lo que implica la existencia de una aceleración centrípeta dirigida siempre hacia el centro del círculo. \\[.35cm]

El MCU es un modelo idealizado que se encuentra frecuentemente en fenómenos naturales y aplicaciones tecnológicas: desde el movimiento de los planetas y satélites artificiales hasta mecanismos rotatorios en maquinaria industrial. \\[.35cm]

En el MCU, aunque un objeto se mueva con rapidez constante, su velocidad cambia constantemente de dirección porque sigue una trayectoria curva.\\[.35cm]

La aceleración centrípeta es la responsable de este cambio de dirección. Se trata de una aceleración que:

\begin{itemize}
    \item siempre apunta hacia el centro del círculo.
  
    \item no cambia la rapidez del objeto, solo la dirección de su velocidad.
   
    \item es esencial para mantener al objeto en movimiento circular.
\end{itemize}

Si denotamos por $v$ la rapidez constante del movimiento y $R$ al radio de la trayectoria circular, la aceleraci\'on centr\'ipeta, denotada $a_\text{rad}$, est\'a dada por
\begin{equation}
a_{\text{rad}}=\frac{v^{2}}{R}.\end{equation}
 (ver Figura \ref{fig:mcu}). \\[.35cm]

El movimiento circular uniforme es resultado de cualquier combinación de fuerzas que produzca una fuerza neta $\sum \vec{F}$
de magnitud constante y dirigida siempre hacia el centro del círculo; por lo tanto, al aplicar la Segunda Ley de Newton a un MCU, tenemos que 
\begin{equation}
    \sum \vec{F}=ma_\text{rad}=m\frac{v^2}{R}
\end{equation}
En el contexto del Movimiento Circular Uniforme, el período se define como el tiempo que tarda un objeto en dar una vuelta completa alrededor de la trayectoria circular. Se representa con la letra $T$ y su unidad en el Sistema Internacional es el segundo (s).\\[.35cm]

El periodo puede determinarse a partir de los par\'ametros del MCU:
\begin{equation}
T=\frac{2\pi R}{v};
\end{equation}
 de manera que la aceleraci\'on centr\'ipeta puede escribirse tambi\'en como
 \begin{equation}
a_{\text{rad}}=\frac{4\pi^{2}R}{T^{2}}
\end{equation}

 
 Al inverso del per\'iodo se le llama \textit{frecuencia}
 \begin{equation}
 f=\frac{1}{T},
 \end{equation}
y representa el n\'umero de revoluciones que completa el m\'ovil por unidad de tiempo. \\[.35cm]

Si el periodo se expresa en segundos, la frecuencia tiene unidades de $\U{s}^{-1}$, lo cual recibe el nombre de \textit{hertz} (Hz):
$$1\U{s}^{-1}=1\U{Hz}.$$


\begin{figure}[ht!]
\centering
\includegraphics[width=5in]{../figuras/mcu.pdf}
\caption{Velocidad y aceleraci\'on de una part\'icula con movimiento circular uniforme.}
\label{fig:mcu}
\end{figure}

\newpage
\section{Movimiento relativo}

\tdplotsetmaincoords{60}{110} 
\begin{figure}[h!]
\begin{center}
\input{planeamiento_didactico/sesiones/movimiento_relativo.tex}
\caption{Movimiento relativo: El observador $A$ se encuentra en reposo, y el observador $B$ se mueve con velocidad constante $\vec{v}_{\text{B/A}}$ respecto de $A$. $B$ observa el objeto $P$ moverse con velocidad $\vec{v}_{\text{P/B}}$}
\label{fig:movrel}
\end{center}
\end{figure}

Cuando un mismo movimiento es descrito por medio de dos sistemas de referencia distintos,  es posible relacionar ambas descripciones si se conoce el \textit{movimiento relativo} de un sistema de referencia respecto al otro. \\[.35cm]
\begin{comment}
    
\begin{figure}[h!]
%\hspace{-5ex}
\centering
\includegraphics[height=1.75in]{../figuras/relativo1-1}
\caption{Ejemplo de \textit{Movimiento relativo}.}
\label{fig:movrel2}
\end{figure}
\end{comment}

Considere dos observadores $A$ y $B$. Por simplicidad, considere que $A$ est\'a en reposo y $B$ se mueve con una velocidad constante $\vec{v}_{\text{B/A}}$ respecto a $A$. Si $B$ observa el movimiento de un objeto $P$ con velocidad constante $\vec{v}_{\text{P/B}}$, $A$ observar\'a a $P$ moverse con una velocidad

\begin{equation}
\vec{v}_{\text{P/A}}= \vec{v}_{\text{P/B}}+\vec{v}_{\text{B/A}};
\label{eq:rel1}
\end{equation}

\begin{figure}[h!]
%\hspace{-5ex}
\centering
\includegraphics[height=1.75in]{planeamiento_didactico/figuras/relativo1.png}
\caption{\textit{Esquema de velocidades} $-$ Movimiento relativo.}
\label{fig:movrel2}
\end{figure}

Equivalentemente, se puede conocer la posici\'on de $P$ respecto a $A$, si se conoce la posici\'on de $P$ respecto a $B$:


\begin{equation}
\vec{r}_{\text{P/A}}= \vec{r}_{\text{P/B}}+\vec{r}_{\text{B/A}}.
\label{eq:rel2}
\end{equation}

Esquem\'aticamente, estas relaciones se muestran en la Figura \ref{fig:movrel}. \\[.35cm]


\textcolor{c1}{\textbf{NOTA:}} Al ser las ecuaciones \eqref{eq:rel1} y \eqref{eq:rel2} expresiones vectoriales; estas se pueden expresar adem\'as como tres ecuaciones escalares; por ejemplo, se tiene que cumplir que
$$x_{\text{P/A}}= x_{\text{P/B}}+x_{\text{B/A}},
 $$
$$y_{\text{P/A}}= y_{\text{P/B}}+y_{\text{B/A}},
 $$ 
$$z_{\text{P/A}}= z_{\text{P/B}}+z_{\text{B/A}}.
 $$
\newpage
\section{Ley de Newton de la Gravitaci\'on}

\tdplotsetmaincoords{-20}{10}
%define polar coordinates for some vector
%TODO: look into using 3d spherical coordinate system
\pgfmathsetmacro{\rvec}{0}
\pgfmathsetmacro{\thetavec}{0}
\pgfmathsetmacro{\phivec}{0}
\begin{figure}[!h]
\begin{tikzpicture}[tdplot_main_coords, auto]
\begin{scope}[rotate=0,scale=1.3]
%\draw [help lines] (-5,-1) grid (5,3);
\begin{scope}[opacity=.8]
\input{eje3d}
\end{scope}
\draw[dotted, very thin] (0,0) circle [x radius=5cm, y radius=5mm, rotate=-10];
\node[opacity=.8] at (5,0) {\includegraphics[width=.15in]{planeamiento_didactico/figuras/tierra2.jpg}};
\node[opacity=.5] at (-3,0) {\includegraphics[width=.35in]{planeamiento_didactico/figuras/sun.png}};
\draw[<->] (-3,.75)--(5,.75);
\node at (0,1) [above]{r};
\node at (-3.5,0) {$m_{1}$};
\node at (5.25,0)[right] {$m_{2}$};
\end{scope}
\end{tikzpicture}
\caption{Interacci\'on gravitacional.}
\label{fig:fgrav}
\end{figure}


La \textit{ley de gravitaci\'on universal} de Newton establece la relaci\'on que caracteriza la interacci\'on entre dos masas ($m_{1}$ y $m_{2}$) separadas cierta distancia ($r$). Esta interacci\'on es una fuerza \textbf{atractiva}. \\[.35cm]

La magnitud de la fuerza gravitacional que ejerce $m_{1}$ sobre $m_{2}$, $\vec{F}_{\text{g}}$, est\'a dada por
\begin{equation}
F_\text{g}=-\frac{Gm_{1}m_{2}}{r^{2}}
\end{equation}
donde  $G$ es una constante de proporcionalidad llamada \textit{constante de gravitaci\'on universal}. El t\'ermino \textit{universal} hace referencia al hecho que se considera constante en todo el universo. En el S.I $$G=6,67\tentimes{-11}\frac{\U{N}\cdot \U{m}^{2}}{\U{kg}^{2}}. $$



La fuerza gravitacional permite explicar tanto el movimiento de los planetas alrededor del Sol como el de los sat\'elites naturales y artificiales alrededor de los planetas (ver Figura \ref{fig:fgrav}). \\[.35cm]

Las fuerzas gravitacionales siempre actúan a lo largo de la línea que une las dos
partículas, y forman un par acción-reacción.




\begin{comment}
\begin{figure}
\centering
\includegraphics[width=3.5in]{planeamiento_didactico/figuras/satelite.png}
\caption{Movimiento de sat\'elites.}
\label{fig:potgrav}
\end{figure}
\end{comment}

\subsubsection{Gravitación y cuerpos esféricamente simétricos}

La interacción gravitacional entre dos cuerpos con distribuciones de masa esféricamente simétricas (ya sean sólidas o huecas) es la misma que si toda
la masa estuviera concentrada en el centro, como se muestra en la Figura \ref{fig:fuerza-g}.\\[.35cm]
\begin{figure}[h!]
    \centering
\includegraphics[width=280px]{planeamiento_didactico/figuras/fuerza-g.pdf}
    \caption{El efecto gravitacional afuera de
cualquier distribución de masa esféricamente simétrica es el mismo que si toda la
masa estuviera concentrada en su centro.}
    \label{fig:fuerza-g}
\end{figure}
Si modelamos la Tierra como un cuerpo esféricamente simétrico de masa $m_\text{T}$, la fuerza
que ejerce sobre una partícula o un cuerpo esféricamente simétrico con masa $m$, a una
distancia $r$ entre los centros, es
\begin{equation}
    F_\text{g}=\frac{Gm_\text{T}m}{r^2}
\end{equation}
siempre y cuando el cuerpo se encuentre en el exterior de la Tierra.

\subsubsection{Peso y gravedad}

Considere un objeto de masa $m$ en la superficie terrestre
$$F_{\text{g}}=\frac{Gm_{\text{T}}m}{R^{2}_{\text{T}}}=mg \R g=\frac{Gm_{\text{T}}}{R^{2}_{\text{T}}}=9,81\U{m/s}^{2},$$
donde $m_{\text{T}}=5,98\tentimes{24}\U{kg}$ es la masa del planeta Tierra y $R_{\text{T}}=6380\U{km}$ es su radio. Conforme el objeto se aleje de la superficie terrestre; su peso ir\'a disminuyendo (ver Figura \ref{fig:potgrav}); de forma tal que 
\begin{equation}
g=g(r)=\frac{Gm_{\text{T}}}{(r-R_\text{T})^{2}},
\end{equation}
donde $r$ es la distancia entre el centro de la Tierra y el objeto.

\begin{figure}[h!]
\includegraphics[width=5in]{planeamiento_didactico/figuras/graficapotencial.png}
\caption{Peso como funci\'on de la distancia.}
\label{fig:potgrav}
\end{figure}